\documentclass[../../../../main.tex]{subfiles}
\begin{document}

\begin{flushright}
\footnotesize\textit{【自動文字起こし・要確認】}
\end{flushright}

{\bf [解]}
$e(\theta)=\cos\theta+i\sin\theta$とおく．

\textbf{(1)} $P_0=0,\ P_1=1+i$．$\overrightarrow{P_{n-1}P_n}$を表す複素数を$d_n$とすると，題意から
\begin{align*}
d_{n+1}=\frac23e\left(\frac\pi3\right)d_n \quad\cdots\text{②}
\end{align*}
$d_1=1+i$（①）とあわせて
\begin{align*}
d_n=\left\{\frac23e\left(\frac\pi3\right)\right\}^{n-1}(1+i)
\end{align*}
だから
\begin{align*}
\overrightarrow{OP_n}=\overrightarrow{OP_0}+d_1+\cdots+d_n=\sum_{k=1}^n\left\{\frac23e\left(\frac\pi3\right)\right\}^{k-1}(1+i)=(1+i)\cdot\frac{1-\left\{\frac23e(\pi/3)\right\}^n}{1-\frac23e(\pi/3)}
\end{align*}
$\left|\dfrac23e(\pi/3)\right|=\dfrac23<1$だから$n\to\infty$で
\begin{align*}
P_\infty=\frac{1+i}{1-\frac23e(\pi/3)}
\end{align*}
$e(\pi/3)=\dfrac12+\dfrac{\sqrt3}2i$より$\dfrac23e(\pi/3)=\dfrac13+\dfrac{\sqrt3}3i$，$1-\dfrac23e(\pi/3)=\dfrac23-\dfrac{\sqrt3}3i=\dfrac{2-\sqrt3i}3$だから
\begin{align*}
P_\infty=\frac{3(1+i)}{2-\sqrt3i}=\frac{3(1+i)(2+\sqrt3i)}{(2-\sqrt3i)(2+\sqrt3i)}=\frac{3(1+i)(2+\sqrt3i)}7
\end{align*}
$(1+i)(2+\sqrt3i)=(2-\sqrt3)+(2+\sqrt3)i$だから
\begin{align*}
P_\infty=\frac37\left\{(2-\sqrt3)+(2+\sqrt3)i\right\}
\end{align*}

\textbf{(2)} $Q_0=0,\ Q_1=z$．(1)と同様に，$\overrightarrow{Q_{n-1}Q_n}$を表す複素数を$\beta_n$とすると$\beta_1=z$，
\begin{align*}
\beta_{n+1}=\frac12e\left(\frac\pi6\right)\beta_n \quad\therefore\ \beta_n=\left\{\frac12e\left(\frac\pi6\right)\right\}^{n-1}z
\end{align*}
\begin{align*}
\overrightarrow{OQ_n}=\sum_{k=1}^n\beta_k=z\cdot\frac{1-\left\{\frac12e(\pi/6)\right\}^n}{1-\frac12e(\pi/6)} \ \xrightarrow{n\to\infty}\ Q_\infty=\frac z{1-\frac12e(\pi/6)}
\end{align*}
$e(\pi/6)=\dfrac{\sqrt3}2+\dfrac12i$より$\dfrac12e(\pi/6)=\dfrac{\sqrt3}4+\dfrac14i$，$1-\dfrac12e(\pi/6)=\dfrac{4-\sqrt3}4-\dfrac14i=\dfrac{(4-\sqrt3)-i}4$だから
\begin{align*}
Q_\infty=\frac{4z}{(4-\sqrt3)-i}
\end{align*}

$Q_\infty=P_\infty$より
\begin{align*}
\frac{4z}{(4-\sqrt3)-i}=\frac37\left\{(2-\sqrt3)+(2+\sqrt3)i\right\}
\end{align*}
\begin{align*}
z=\frac3{28}\left\{(2-\sqrt3)+(2+\sqrt3)i\right\}\left\{(4-\sqrt3)-i\right\}
\end{align*}
右辺の$\{\ \}\{\ \}$を展開する（実部$=pr-qs$，虚部$=ps+qr$，ただし$p=2-\sqrt3,q=2+\sqrt3,r=4-\sqrt3,s=-1$として）：
\begin{align*}
pr=(2-\sqrt3)(4-\sqrt3)=11-6\sqrt3,\qquad qs=-(2+\sqrt3)
\end{align*}
\begin{align*}
pr-qs=11-6\sqrt3+2+\sqrt3=13-5\sqrt3
\end{align*}
\begin{align*}
ps=-(2-\sqrt3),\qquad qr=(2+\sqrt3)(4-\sqrt3)=5+2\sqrt3
\end{align*}
\begin{align*}
ps+qr=-(2-\sqrt3)+5+2\sqrt3=3+3\sqrt3=3(1+\sqrt3)
\end{align*}
よって
\begin{align*}
z=\frac3{28}\left\{(13-5\sqrt3)+3(1+\sqrt3)i\right\}
\end{align*}


\end{document}
